% quark_scan_assumptions_compact.tex
% Self-contained walkthrough of the quark-sector scan pipeline.
% This compact note is the tracked quark-scan writeup in the repo.
\documentclass[11pt]{article}

\usepackage[margin=1in]{geometry}
\usepackage{parskip}
\usepackage{amsmath,amssymb}
\usepackage{booktabs}
\usepackage{xcolor}
\usepackage{hyperref}
\usepackage{enumitem}

\hypersetup{colorlinks=true,linkcolor=blue,citecolor=blue,urlcolor=blue}

\newcommand{\Mkk}{M_{\mathrm{KK}}}
\newcommand{\LamIR}{\Lambda_{\mathrm{IR}}}
\newcommand{\Mpl}{M_{\mathrm{Pl}}}
\newcommand{\gs}{g_s^*}

\title{Quark Sector Scan: Apr 16}
\author{}
\date{}

\begin{document}
\maketitle

%% ====================================================================
\section{Geometry}

We use the following metric convention
\begin{equation}
  ds^2 = (kz)^{-2}\bigl(\eta_{\mu\nu}\,dx^\mu\,dx^\nu + dz^2\bigr),
  \qquad \eta = \mathrm{diag}(-1,+1,+1,+1),
\end{equation}
with the UV brane at $z = 1/k$ and the IR brane at $z = 1/\LamIR$.

Two inputs fix everything: the AdS curvature $k$ and the IR scale $\LamIR$.
We set $k = \Mpl = 1.2209\times10^{19}$~GeV throughout and scan over $\LamIR$
from 500~GeV to 20~TeV. The warp factor is $\varepsilon = \LamIR / k \sim
10^{-15}$. The Higgs is taken to be localized on the IR brane, and we use
$v = 174$~GeV as the electroweak VEV (i.e.\ $\langle H \rangle =
v_{\mathrm{SM}}/\sqrt{2} = 174$~GeV). Backreaction is neglected.


%% ====================================================================
\section{Zero-mode profiles}

Each bulk fermion has a dimensionless mass parameter $c$, which
controls where its zero mode is localized in the extra dimension. 
\textcolor{red}{\textbf{Possible to do:} We can make this the full integral rather than just the IR overlap}
The key
quantity is the zero-mode's overlap with the IR brane, since the Higgs
(and therefore the Yukawa interaction) lives there:
\begin{equation}
  f_{\mathrm{IR}}^2(c,\varepsilon)
    = \frac{\tfrac{1}{2} - c}{1 - \varepsilon^{1-2c}}\,.
\end{equation}
This is the point of the RS flavor story: the mass hierarchy comes from the
exponential sensitivity of $f_{\mathrm{IR}}$ to $c$, not from hierarchical
5D Yukawa couplings.


%% ====================================================================
\section{Yukawa spurion construction}

The 5D Yukawa matrices $Y_u$ and $Y_d$ are complex $3\times3$ matrices. We
parametrize each via a singular value decomposition:
\begin{equation}
  Y = \texttt{overall\_scale} \times U_{\mathrm{left}}
    \cdot \operatorname{diag}(\sigma_1,\sigma_2,\sigma_3)
    \cdot U_{\mathrm{right}}^\dagger,
\end{equation}
where $\sigma_i > 0$ are ordered light-to-heavy and $U_{\mathrm{left}}$ is a unitary parametrized by three angles $(\theta_{12}, \theta_{13},
\theta_{23})$ and a CP phase $\delta$. The prefactor \texttt{overall\_scale}
is a common real positive number multiplying both $Y_u$ and $Y_d$; it controls
the overall magnitude of the Yukawa entries.

We freeze $U_{\mathrm{right}} = I$
for both sectors, which is a restriction: it means $Y^\dagger Y =
\operatorname{diag}(\sigma^2)$ is automatically diagonal, while
$Y Y^\dagger = U_{\mathrm{left}} \operatorname{diag}(\sigma^2)
U_{\mathrm{left}}^\dagger$ is generally dense. Before the fit starts,
\texttt{overall\_scale} is absorbed into the singular values and set to~1.

Setting $U_{\mathrm{right}} = I$ matters because the right-handed
bulk mass matrices $C_u = Y_u^\dagger Y_u$ and $C_d = Y_d^\dagger Y_d$ are
diagonal by construction, so the dominant source of flavor misalignment is the
left-handed matrix $C_Q$. (Right-handed off-diagonal couplings can still appear
after the SVD rotation to the mass basis, but the restriction to diagonal $C_u$,
$C_d$ removes much of the freedom that a general scan would have.)
This restricts the scan to 15 free parameters (6 singular values +
8 left-rotation angles + 1 overall scale) out of the 36 that a pair of
general $3\times3$ complex matrices would have.



%% ====================================================================
\section{Bulk mass matrices and the role of $r$}

From the Yukawa spurions we build three Hermitian matrices that determine the
bulk localization of the quark zero modes:
\begin{equation}
  C_u = Y_u^\dagger Y_u, \qquad
  C_d = Y_d^\dagger Y_d, \qquad
  C_Q = r\,(Y_u Y_u^\dagger) + (Y_d Y_d^\dagger).
\end{equation}

In the NMFV literature the full form is $C_Q = \alpha_Q\,\mathbf{I} +
\beta_Q\,Y_d Y_d^\dagger + \gamma_Q\,Y_u Y_u^\dagger$, where $\alpha_Q$ is
a flavor-universal bulk mass and $r = \gamma_Q / \beta_Q$. This initial scan omits the
$\alpha_Q\,\mathbf{I}$ term entirely. 
\textcolor{red}{\textbf{To do}: rerun the scan to include $\alpha_Q$. This is very simple to do in the code. We can easily change the assumptions and test new models quickly!} clamps the eigenvalues into a bounded window after
diagonalization. This is not equivalent to including $\alpha_Q$: it preserves
the eigenvectors but distorts the eigenvalue splittings in a way that has no
counterpart in the MFV spurion expansion.

The parameter $r$ controls alignment. When $r = 0$, $C_Q \propto Y_d Y_d^\dagger$,
which can be diagonalized simultaneously with $Y_d$. This follows because $U_{\mathrm{right}} = I$,
so $Y_d = U_{\mathrm{left}}\,\mathrm{diag}(\sigma)$ and the
rotation that diagonalizes $C_Q = Y_d Y_d^\dagger$ also diagonalizes
$Y_d$ itself. In that basis the
left handed overlap matrix $F_Q$ is diagonal, so there are no tree level FCNCs
in the down sector. As $r$ increases, the $Y_u Y_u^\dagger$ term introduces off diagonal entries
in $C_Q$, misaligning the left handed quarks from the down mass eigenbasis and
generating tree level FCNCs.


%% ====================================================================
\section{From spurion eigenvalues to physical bulk masses}

We diagonalize each $C_X$ to get eigenvalues $\lambda_i$ and a rotation
matrix. The question is how to turn those eigenvalues into physical bulk mass
parameters $c_i$. The two frameworks we want to be compatible with handle
this differently:
\begin{itemize}
  \item In 0710.1869, specific benchmark $c$-values are just quoted directly
    (Table~I), but with ``universal terms and overall order one coefficients
    omitted for simplicity.'' The bulk masses are input data. This is what I was telling you about earlier.
  \item In the Csaki et al.\ (0907.0474) convention, the full
    $C_Q = \alpha_Q\,\mathbf{I} + \beta_Q\,Y_d Y_d^\dagger +
    \gamma_Q\,Y_u Y_u^\dagger$ is kept, and the eigenvalues of $C_Q$ are
    directly the bulk masses. The $\alpha_Q$ term keeps the spectrum in a
    physical window.
\end{itemize}

Since this initial scan omits $\alpha_Q$, the eigenvalues of $C_Q$ are raw
non-negative numbers on $[0,\infty)$ with nothing anchoring them in a
physical range. To keep all bulk masses inside $[c_{\mathrm{ir}},
c_{\mathrm{uv}}]$ we apply the map
\begin{equation}
  c = c_{\mathrm{uv}} - (c_{\mathrm{uv}} - c_{\mathrm{ir}})
    \;\frac{\lambda}{\lambda + 1}\,,
\end{equation}
with $c_{\mathrm{uv}} = 0.72$, $c_{\mathrm{ir}} = 0.30$, giving
$c = 0.72 - 0.42\,\lambda/(\lambda+1)$. This maps $\lambda = 0 \to c = 0.72$
(UV-localized), $\lambda \to \infty \to c \to 0.30$ (IR-localized), and is
applied identically to all three sectors ($Q$, $u$, $d$).


The choice of $\lambda/(\lambda+1)$ is not arbitrary; Expanding at small $\lambda$,
\begin{equation}
  c \;\approx\; 0.72 - 0.42\,\lambda \qquad (\lambda \ll 1),
\end{equation}
which is exactly the linear identification $c = \alpha - \beta\,\lambda$
that the Csaki et al.\ parametrization would give. Over the actual scan,
the $C_Q$, $C_u$, $C_d$ eigenvalues have median $\lambda \approx 0.02$,
and in that regime the sigmoid and the affine map agree to better than
$10^{-3}$. The sigmoid only departs from linearity in the large-$\lambda$
tail (eigenvalues up to $\sim\!11$), where you end up having that the affine map would push
$\sim\!26\%$ of points below $c_{\mathrm{ir}} = 0.30$ into unphysical
territory.


The window $[0.30, 0.72]$ is a design choice: $c = 0.30$ gives
$f_{\mathrm{IR}} \sim 0.45$, which forces the top Yukawa to $\sim 2.4$
to reproduce $m_t = 172$~GeV (perturbative but above the naive $O(1)$
expectation). Using the same map for all three sectors removes the
independent baselines ($\alpha_Q$, $\alpha_u$, $\alpha_d$) that the full
parametrization would provide.

\textcolor{red}{\textbf{Note:} With the explicit $(\alpha_Q, \beta_Q,
\gamma_Q)$ parametrization the sigmoid goes away entirely. That is, eigenvalues are
directly the bulk masses. The sigmoid
is an assumption that is likely to bias the top sector masses and the FCNC
bounds.}

\textcolor{red}{\textbf{Note:} The entire analysis is tree-level.
In the Csaki et al.\ (0907.0474) ``shining'' framework, up type Yukawa
loops induce radiative flavor misalignment, generating off-diagonal
down type Yukawas and kinetic mixings that can be numerically important
near current bounds.}

%% ====================================================================
\section{Zero-mode overlaps, basis rotation, and mass matrices}

With the physical bulk masses in hand, we compute the IR-brane overlaps
$F_Q[i] = f_{\mathrm{IR}}(c_Q[i], \varepsilon)$ for the left-handed doublets,
and similarly $F_u$, $F_d$ for the right-handed singlets. We then rotate the
Yukawas into the bulk-mass eigenbasis, and the effective 4D quark mass matrices are then
\begin{equation}
  M_{u,d} = 2v \cdot \operatorname{diag}(F_Q) \cdot Y_{u,d}^{\mathrm{bulk}}
    \cdot \operatorname{diag}(F_{u,d}),
\end{equation}
where $2v = 348$~GeV $\approx \sqrt{2}\,v_{\mathrm{SM}}$. The 5D Yukawas
in this convention are dimensionless. 


%% ====================================================================
\section{Physical masses, CKM, and the fit}

Each mass matrix is decomposed by SVD:
$M = U_L \operatorname{diag}(m_1, m_2, m_3) U_R^\dagger$, giving the physical
quark masses as singular values. The CKM matrix is $V_{\mathrm{CKM}} =
(U_L^u)^\dagger U_L^d$, from which we extract $|V_{us}|$, $|V_{cb}|$,
$|V_{ub}|$, and the Jarlskog invariant $J$.

We wrap the whole chain (Yukawa parameters $\to$ masses + CKM) in an optimizer that adjusts 14 free parameters (6
log-singular-values $+$ 8 left-rotation angles) to match the observed quark
masses and CKM elements. The residual is a 10-component vector: 6 log-mass
ratios plus 4 fractional CKM deviations. We define the fit score as the RMS
of these 10 residuals; points with score $> 0.1$ (roughly $> 10\%$ average
deviation) are discarded.

The target masses are PDG~2024 $\overline{\mathrm{MS}}$ values RG-evolved
to a single common scale $\mu_{\mathrm{common}} = m_t(m_t) = 163.5\,\mathrm{GeV}$
(this is also \texttt{qcd.constants.M\_TOP\_MS} in the codebase). The
running uses the standard four-loop $\overline{\mathrm{MS}}$ mass anomalous
dimension with three-loop threshold matching at $m_b$ and (where needed)
$m_c$ following Chetyrkin--Kniehl--Steinhauser. No path runs below $m_c$.

The Wilson-coefficient $\alpha_s$ reference scale is independent of the
mass-target scale and stays at $\mu = 3\,\mathrm{TeV}$ for the
$\Delta F = 2$ matching; the two scales are kept orthogonal on purpose.
The CKM matrix is held constant between $\mu_{\mathrm{common}}$ and the
hadronic matching scale $\mu_{\mathrm{had}} \ge m_W$; this is the
conventional one-loop QCD-invariant treatment (PDG~2024 §12).

Acceptance gating uses per-quark and per-CKM-element $2\sigma$ relative
tolerances. The mass tolerances are derived from the PDG~2024 quoted
$1\sigma$ uncertainties (with a deterministic floor of $0.003$ to absorb
optimizer precision). The CKM gate uses
\[
\{0.0044,\, 0.072,\, 0.10,\, 0.085\}
\]
for
$\{|V_{us}|,\, |V_{cb}|,\, |V_{ub}|,\, J\}$. The empirical
floor will be tightened to the 95th-percentile per-quark log-residual
once the user runs the Phase-0 calibration sweep
(\texttt{scripts/calibrate\_phase0.py}); until then, the deterministic
floor of $0.003$ is in force.

\paragraph{Plan-B audit gate.}
If the production scan returns zero accepted points, the fallback policy
relaxes tolerances only after a per-quark median-residual audit
(committed to \texttt{data/per\_quark\_median\_residual\_pdg2024.json})
and a written physics-tension justification in this document. No quark
whose median residual already sits at or above $2\sigma$ may be relaxed
to $3\sigma$ without that justification. See plan v3 §9.2 for the
decision tree.

\paragraph{Re-running the scan.}
The new accepted-points artifact is generated by
\begin{verbatim}
python scripts/export_accepted_quark_scan_with_yukawas.py \
       --output data/accepted_points_with_yukawas_pdg2024.csv
\end{verbatim}
and accompanied by a fresh \texttt{.provenance.json} containing the PDG
edition, target label \texttt{pdg-2024-msbar-mu-mt-v1}, mass-target
scale $\mu = 163.5$~GeV, WC reference scale $\mu = 3\,\mathrm{TeV}$,
git commit, and the per-quark / per-CKM tolerance vectors used.


%% ====================================================================
\section{The scan}

The scan sweeps three parameters on a grid:

\begin{center}
\begin{tabular}{@{}llll@{}}
\toprule
Parameter & Grid & Range & Points \\
\midrule
$r$ & log-uniform & 0.02--2.0 & 100 \\
\texttt{overall\_scale} & linear & 1.5--6.0 & 20 \\
$\LamIR$ & log-uniform & 500--20\,000~GeV & 50 \\
\bottomrule
\end{tabular}
\end{center}

$r$ controls the degree of alignment in $C_Q$ (previous sections).
\texttt{overall\_scale} sets the overall Yukawa magnitude. Different values
lead to different internal Yukawa structures even after the fitter re-optimizes
to hit the same quark masses: if \texttt{overall\_scale} is large, all Yukawa
entries are large and the fitter has to make the singular values small to
compensate; if it is small, the singular values have to be large. These
different configurations produce different off-diagonal entries in the mass
basis and therefore different FCNC predictions, even though the diagonal masses
come out the same. Since the fitter freely adjusts all
six singular values, \texttt{overall\_scale} is formally reabsorbable and is
not an independent physical degree of freedom. Its role is just that of a scan hyperparameter, since different values seed the optimizer in
different basins of the fit landscape.


Over the scan range 1.5 to 6.0, the largest singular value
(top-like) runs from $\sim 1.5$ to $\sim 6$ and the smallest (up-like) from
$\sim 0.005$ to $\sim 0.02$. At the high end the top Yukawa is well above the
naive $O(1)$ expectation. $\LamIR$ sets the KK scale and therefore the
$1/\Mkk^2$ suppression of all Wilson coefficients.

Total: $100 \times 20 \times 50 = 100{,}000$ points, each fitted independently
from the same deterministic seed.

At scan time several quantities are held fixed. The KK-mass convention is
$\xi_{\mathrm{KK}} = 1.0$, meaning $\Mkk = \LamIR$ in the Wilson coefficient
propagator; the enhanced
KK-gluon coupling is $\gs = 3.0$.

\textcolor{red}{\textbf{To do:} Future scans should set
$\xi_{\mathrm{KK}}$ at the physical first KK-gluon mass from the
gauge-sector Bessel root) so that $\Mkk = m_g^{(1)}$ from the start.}

\textcolor{red}{\textbf{To do:} Confirm $\gs$ with Lisa. Tree-level
$\gs \sim 6$ vs.\ one-loop $\gs \sim 3$ is a factor of 4 in ${\gs}^2$.}


%% ====================================================================
\section{KK-gluon couplings in the mass basis}

After the fit converges we have the diagonal overlaps $F_Q$, $F_u$, $F_d$ and
the left/right rotation matrices $U_L^{u,d}$, $U_R^{u,d}$ from the SVD. We
model the KK-gluon coupling to zero-mode quarks as proportional to
$\operatorname{diag}(F^2)$ in the bulk-mass eigenbasis.
In the physical mass basis this picks up off-diagonal entries from the
rotations:
\begin{equation}
  g_{L}^{\mathrm{down}}[i,j] = \gs \cdot
    \bigl(U_L^{d\,\dagger} \operatorname{diag}(F_Q^2)\, U_L^d\bigr)_{ij},
\end{equation}
and analogously for $g_R^{\mathrm{down}}$ (using $U_R^d$, $F_d$),
$g_L^{\mathrm{up}}$ (using $U_L^u$, $F_Q$), and $g_R^{\mathrm{up}}$ (using
$U_R^u$, $F_u$). Note that $F_Q$ enters both LH couplings because $u_L$ and
$d_L$ share the same $SU(2)$ doublet.


%% ====================================================================
\section{Wilson coefficients}

Integrating out the KK-gluon at tree level at the scale $\mu = \Mkk$ gives
four $\Delta F = 2$ operators. The operators below are labeled
$(Q_4^{\mathrm{LR}}, Q_5^{\mathrm{LR}})$, which is the SUSY basis
convention (Gabbiani et al.), not the BMU basis $(Q_1^{\mathrm{LR}},
Q_2^{\mathrm{LR}})$. The anomalous dimension matrix in
Section~10 and the matrix elements in Section~12 are also written in
the SUSY basis. 
\begin{align}
  C_1^{\mathrm{VLL}} &= \frac{(g_L[i,j])^2}{6\,\Mkk^2}\,, \qquad
  C_1^{\mathrm{VRR}} = \frac{(g_R[i,j])^2}{6\,\Mkk^2}\,, \\[4pt]
  C_4^{\mathrm{LR}} &= -\frac{g_L[i,j]\,g_R[i,j]}{\Mkk^2}\,, \qquad
  C_5^{\mathrm{LR}} = \frac{g_L[i,j]\,g_R[i,j]}{3\,\Mkk^2}\,.
\end{align}
The five meson systems and their quark-pair indices are: $\epsilon_K$ and
$\Delta M_K$ use $(d,s)$; $B_d$ uses $(d,b)$; $B_s$ uses $(s,b)$; $D^0$
uses $(u,c)$. We include only the first KK-gluon mode (\textcolor{red}{Possible To do: we can include higher modes for the subleading correction}) and use $\Mkk = \LamIR$ as both the
propagator mass and matching scale.


%% ====================================================================
\section{QCD running}

The Wilson coefficients are matched at $\Mkk$ and the hadronic
matrix elements are evaluated at $\mu_{\mathrm{had}} = 2$~GeV, so we evolve
the coefficients down using leading log QCD RG equations. The LR
operators $C_4$ and $C_5$ mix via the $2\times2$ anomalous dimension matrix
\begin{equation}
  \gamma_{\mathrm{LR}} = \begin{pmatrix} 8 & -4 \\ -16/3 & -28/3 \end{pmatrix}.
\end{equation}
This is the anomalous dimension matrix
in the SUSY $(C_4, C_5)$ basis, not the BMU $(C_1^{\mathrm{LR}},
C_2^{\mathrm{LR}})$ basis, which matches the code. 
The dominant eigenvalue is large and positive, giving a $3$ to $5\times$
enhancement of the LR operators at low scales. Combined with the chiral
enhancement factor $r_\chi = (m_K/(m_s + m_d))^2 \approx 25$ in the kaon
matrix elements, this is why $\epsilon_K$ is the most powerful constraint.

The running crosses one flavor threshold at $m_b = 4.18$~GeV ($n_f = 5 \to 4$);
$m_c = 1.27$~GeV is below $\mu_{\mathrm{had}}$. We use one-loop $\alpha_s$
running with $\alpha_s(M_Z) = 0.1179$ and continuous threshold matching.


%% ====================================================================
\section{Hadronic matrix elements and observables}

The evolved Wilson coefficients at $\mu_{\mathrm{had}} = 2$~GeV are
contracted with hadronic matrix elements to get the NP mixing amplitude:
\begin{equation}
  M_{12}^{\mathrm{NP}} = C_1^{\mathrm{VLL}} \cdot \mathrm{ME}_{\mathrm{VLL}}
    + C_1^{\mathrm{VRR}} \cdot \mathrm{ME}_{\mathrm{VLL}}
    + C_4^{\mathrm{LR}} \cdot \mathrm{ME}_{\mathrm{LR4}}
    + C_5^{\mathrm{LR}} \cdot \mathrm{ME}_{\mathrm{LR5}},
\end{equation}
where VRR and VLL share the same matrix element by parity. 

\textcolor{red}{\textbf{To do:} I am pretty sure I have a factor of 2 subtlety here I need to fix. The matrix
elements are
\begin{equation}
  \mathrm{ME}_{\mathrm{VLL}} = \tfrac{2}{3} f_P^2 m_P B_1, \qquad
  \mathrm{ME}_{\mathrm{LR4}} = \bigl(\tfrac{r_\chi}{6} + \tfrac{1}{4}\bigr)
    f_P^2 m_P B_4, \qquad
  \mathrm{ME}_{\mathrm{LR5}} = \bigl(\tfrac{r_\chi}{2} + \tfrac{1}{12}\bigr)
    f_P^2 m_P B_5,
\end{equation}
with $r_\chi = (m_P / (m_{q_1} + m_{q_2}))^2$. The matrix elements above
already absorb the $1/(2m_P)$ state-normalization factor, so $M_{12} =
\sum_i C_i \cdot \mathrm{ME}_i$ is the standard off-diagonal mass-matrix
element in GeV and $\Delta m = 2|M_{12}|$. The code computes
$\mathrm{ME}_{\mathrm{VLL}} = \tfrac{2}{3}\,f_P^2\,m_P\,B_1$ and contracts
it directly with $C_i$ to obtain $M_{12}$, but elsewhere we compute
$\langle Q_1\rangle = \tfrac{8}{3}\,f_P^2\,m_P^2\,B_K$ and divides by
$2m_P$ explicitly. With the same $B_1 = B_K$.}


Hadronic inputs (decay constants, bag parameters, quark masses) come from
FLAG~2024 and PDG~2024, with the exception of the $D^0$ LR bag parameters
$B_4$, $B_5$ which are estimated at 1.0.


%% ====================================================================
\section{Acceptance criteria}

For each parameter point we compute the KK-gluon contribution to five meson
mixing observables: $\epsilon_K$, $\Delta M_K$, $\Delta M_{B_d}$,
$\Delta M_{B_s}$, $\Delta M_{D^0}$. Each one gets a ``budget,'' the
amount of room the experimental measurement leaves for new physics beyond the
SM. We form the ratio
\begin{equation}
  \text{ratio} = \frac{\text{KK-gluon prediction}}{\text{budget}}\,,
\end{equation}
and a point passes only if all five ratios are $\le 1$
The budgets are deliberately generous, so these bounds should be read as a
weak ``best-case envelope,'' not a sharp exclusion.

For $\epsilon_K$, the observable is
$\epsilon_K^{\mathrm{NP}} = |\kappa_\epsilon / (\sqrt{2}\,\Delta m_K)| \cdot
|\mathrm{Im}(M_{12}^{\mathrm{NP}})|$
with $\kappa_\epsilon = 0.94$, and the budget is the exp--SM gap:
$|2.228\times10^{-3} - 1.81\times10^{-3}| = 4.18\times10^{-4}$.
The absolute value discards the sign of the so there is no constructive/destructive
interference with SM.

For $\Delta M_K$, the budget is $\Delta m_K^{\mathrm{exp}} / 2$ (the SM
prediction is long distance dominated and unreliable, so the entire
experimental value is treated as room for new physics).

For $B_d$ and $B_s$, the budget is
$\max(\Delta m^{\mathrm{exp}}/2,\; |\Delta m^{\mathrm{exp}} -
\Delta m^{\mathrm{SM}}|/2)$, which in practice is dominated by
$\Delta m^{\mathrm{exp}}/2$. For $D^0$ it is simply
$\Delta m^{\mathrm{exp}}/2$.


\textcolor{red}{\textbf{To do:} I think there is a lot to improve here. Right now we only check whether the RS
contribution fits inside this defined budget for each system independently.}

%% ====================================================================
\section{Post-processing and figures}


Because the scan used $\Mkk = \LamIR$ as the propagator mass, but the
physical first KK-gluon mass is $m_g^{(1)} = \xi_{\mathrm{kk}} \LamIR$ with
$\xi_{\mathrm{kk}} = 2.4487$ (the first gauge-sector Bessel root), the
plotting script relabels the axis to $m_g^{(1)}$ and divides all ratios by
$\xi_{\mathrm{kk}}^2$ to correct the propagator. This does not correct the
QCD running start scale (see the TODO above).

At each $(r, \Mkk)$ grid point there are up to 20 \texttt{overall\_scale}
values. The plotter picks the single \texttt{overall\_scale} that minimizes
$\max(\text{all 5 ratios})$, so that all five contour lines in the figure come
from one physically consistent MFV point.

Both figures are best-case envelopes: the fitter picks the most favorable
Yukawa structure at each grid point, and the plotter picks the most favorable
\texttt{overall\_scale}. So the contour shows the lowest $\Mkk$ at which
any fitted point survives; it is not a statement about
typical NMFV behavior.

\textcolor{red}{\textbf{To do:} These figures answer ``can NMFV ever allow
$\Mkk$ this low?'' Your note asks for an ensemble analysis: random $O(1)$
Yukawas in the explicit $(\alpha_Q, \beta_Q, \gamma_Q)$ parametrization, keep
points reproducing quark masses, histogram $\Mkk$ where
$\mathrm{Im}(C_4^K)$ meets its bound. Both analyses are needed.}

\end{document}
